\documentclass[10pt]{extarticle}
\usepackage[a4paper,margin=1.5cm]{geometry}
\usepackage[T1]{fontenc}
\usepackage{amsmath,amssymb}
\usepackage{graphicx}
\usepackage[export]{adjustbox}
\usepackage{multicol}
\usepackage{enumitem}
\usepackage{booktabs,array,tabularx}
\setlength{\columnsep}{0.7cm}
\setlength{\parindent}{0pt}
\setlength{\parskip}{3pt}
\pagestyle{empty}
\makeatletter
\newcommand{\qfigdef}[4]{\@namedef{qf@file@#1}{#2}\@namedef{qf@w@#1}{#3}\@namedef{qf@wide@#1}{#4}}
\newcommand{\qf@check}[1]{\@ifundefined{qf@file@#1}{\PackageError{score2}{Unknown figure id #1}{}}{}}
\newcommand{\qf@img}[4][]{\includegraphics[width=#2\dimexpr\@nameuse{qf@w@#4}pt\relax,max width=#3,max height=0.45\textheight,#1]{\@nameuse{qf@file@#4}}}
\newif\ifqf@full
\newif\ifqf@boxed
\newsavebox{\qf@box}
\newcommand{\qf@open}{\global\qf@fulltrue\par\global\qf@boxedtrue
  \setbox\qf@box\vbox\bgroup\hsize\textwidth\linewidth\textwidth}
\newcommand{\qf@emit}[1]{\par\egroup\global\qf@boxedfalse
  \ifdim\dimexpr\ht\qf@box+\dp\qf@box\relax<#1\textheight\noindent\usebox{\qf@box}\par\else\unvbox\qf@box\fi}
\newcommand{\qfullsplit}{\ifqf@boxed\qf@emit{0.85}\fi}
\newcommand{\qf@close}{\ifqf@boxed\qf@emit{0.45}\fi\global\qf@fullfalse}
\newcommand{\qfullbegin}{\end{multicols}\qf@open}
\newcommand{\qfullend}{\qf@close\begin{multicols}{2}}
\newcommand{\qfullstart}{\qf@open}
\newcommand{\qfulldone}{\qf@close}
\newcommand{\qfig}[2][1]{\qf@check{#2}%
  \ifqf@full\qf@blk{#1}{#2}\else
  \ifnum\@nameuse{qf@wide@#2}=1 \expandafter\@firstoftwo\else\expandafter\@secondoftwo\fi
  {\end{multicols}\begin{center}\qf@img{#1}{\textwidth}{#2}\end{center}\begin{multicols}{2}}%
  {\qf@blk{#1}{#2}}\fi}
\newcommand{\qf@blk}[2]{\par\smallskip{\centering\qf@img{#1}{\linewidth}{#2}\par}\smallskip}
\newcommand{\qfigi}[2][1]{\qf@check{#2}\qf@img[valign=c]{#1}{\linewidth}{#2}}
\makeatother
\newcommand{\papersec}[1]{\par\bigskip{\centering\Large\bfseries PAPER-#1\par}\medskip}
\newcommand{\partsec}[2]{\par\medskip\noindent{\large\bfseries PART-#1: #2}\par\smallskip}
\newcommand{\sect}[1]{\par\smallskip\noindent\textbf{#1}\par}
\newcommand{\qn}[1]{\par\smallskip\noindent\textbf{#1.}~}
\newcommand{\paraq}[2]{\par\smallskip\noindent\textbf{Paragraph for Questions #1 and #2}\par}
\newenvironment{qoptions}{\begin{itemize}[leftmargin=1.8em,itemsep=1pt,topsep=2pt,label={}]}{\end{itemize}}
\newcommand{\opt}[2]{\item[(#1)] #2}
% figure registry (generated by score2_tools.py finalize)
\qfigdef{p1_phy_s1_q01_1}{JEE Advanced (2021 Series) - Test 03 (01 Mar 2021)-img/p1_phy_s1_q01_1.png}{109.1}{0}
\qfigdef{p1_phy_s1_q03_1}{JEE Advanced (2021 Series) - Test 03 (01 Mar 2021)-img/p1_phy_s1_q03_1.png}{100.4}{0}
\qfigdef{p1_phy_s2_q05_1}{JEE Advanced (2021 Series) - Test 03 (01 Mar 2021)-img/p1_phy_s2_q05_1.png}{95.9}{0}
\qfigdef{p1_phy_s2_q06_1}{JEE Advanced (2021 Series) - Test 03 (01 Mar 2021)-img/p1_phy_s2_q06_1.png}{101.5}{0}
\qfigdef{p1_phy_s2_q07_1}{JEE Advanced (2021 Series) - Test 03 (01 Mar 2021)-img/p1_phy_s2_q07_1.png}{171.1}{0}
\qfigdef{p1_phy_s2_q08_2}{JEE Advanced (2021 Series) - Test 03 (01 Mar 2021)-img/p1_phy_s2_q08_2.png}{153.8}{0}
\qfigdef{p1_phy_s2_q08_1}{JEE Advanced (2021 Series) - Test 03 (01 Mar 2021)-img/p1_phy_s2_q08_1.png}{148.4}{0}
\qfigdef{p1_phy_s2_q09_1}{JEE Advanced (2021 Series) - Test 03 (01 Mar 2021)-img/p1_phy_s2_q09_1.png}{87.7}{0}
\qfigdef{p1_phy_s2_q10_1}{JEE Advanced (2021 Series) - Test 03 (01 Mar 2021)-img/p1_phy_s2_q10_1.png}{150.4}{0}
\qfigdef{p1_phy_s2_q12_1}{JEE Advanced (2021 Series) - Test 03 (01 Mar 2021)-img/p1_phy_s2_q12_1.png}{184.4}{0}
\qfigdef{p1_phy_s3_q01_1}{JEE Advanced (2021 Series) - Test 03 (01 Mar 2021)-img/p1_phy_s3_q01_1.png}{334.5}{1}
\qfigdef{p1_phy_s3_q04_1}{JEE Advanced (2021 Series) - Test 03 (01 Mar 2021)-img/p1_phy_s3_q04_1.png}{110.0}{0}
\qfigdef{p1_phy_s3_q06_1}{JEE Advanced (2021 Series) - Test 03 (01 Mar 2021)-img/p1_phy_s3_q06_1.png}{200.3}{0}
\qfigdef{p1_chem_s2_q05_1}{JEE Advanced (2021 Series) - Test 03 (01 Mar 2021)-img/p1_chem_s2_q05_1.png}{32.3}{0}
\qfigdef{p1_chem_s2_q05_2}{JEE Advanced (2021 Series) - Test 03 (01 Mar 2021)-img/p1_chem_s2_q05_2.png}{32.3}{0}
\qfigdef{p1_chem_s2_q05_D}{JEE Advanced (2021 Series) - Test 03 (01 Mar 2021)-img/p1_chem_s2_q05_D.png}{50.0}{0}
\qfigdef{p1_chem_s2_q06_2}{JEE Advanced (2021 Series) - Test 03 (01 Mar 2021)-img/p1_chem_s2_q06_2.png}{99.9}{0}
\qfigdef{p1_chem_s2_q06_1}{JEE Advanced (2021 Series) - Test 03 (01 Mar 2021)-img/p1_chem_s2_q06_1.png}{178.3}{0}
\qfigdef{p1_chem_s2_q08_1}{JEE Advanced (2021 Series) - Test 03 (01 Mar 2021)-img/p1_chem_s2_q08_1.png}{127.6}{0}
\qfigdef{p1_chem_s2_q12_D}{JEE Advanced (2021 Series) - Test 03 (01 Mar 2021)-img/p1_chem_s2_q12_D.png}{65.3}{0}
\qfigdef{p1_chem_s2_q12_3}{JEE Advanced (2021 Series) - Test 03 (01 Mar 2021)-img/p1_chem_s2_q12_3.png}{46.5}{0}
\qfigdef{p1_chem_s3_q04_1}{JEE Advanced (2021 Series) - Test 03 (01 Mar 2021)-img/p1_chem_s3_q04_1.png}{80.9}{0}
\qfigdef{p1_chem_s3_q04_2}{JEE Advanced (2021 Series) - Test 03 (01 Mar 2021)-img/p1_chem_s3_q04_2.png}{34.3}{0}
\qfigdef{p1_chem_s3_q04_3}{JEE Advanced (2021 Series) - Test 03 (01 Mar 2021)-img/p1_chem_s3_q04_3.png}{55.0}{0}
\qfigdef{p1_chem_s3_q04_4}{JEE Advanced (2021 Series) - Test 03 (01 Mar 2021)-img/p1_chem_s3_q04_4.png}{62.6}{0}
\qfigdef{p1_chem_s3_q04_5}{JEE Advanced (2021 Series) - Test 03 (01 Mar 2021)-img/p1_chem_s3_q04_5.png}{80.3}{0}
\qfigdef{p1_chem_s3_q04_6}{JEE Advanced (2021 Series) - Test 03 (01 Mar 2021)-img/p1_chem_s3_q04_6.png}{30.4}{0}
\qfigdef{p1_math_s2_q11_1}{JEE Advanced (2021 Series) - Test 03 (01 Mar 2021)-img/p1_math_s2_q11_1.png}{98.4}{0}
\qfigdef{p2_phy_s1_q01_1}{JEE Advanced (2021 Series) - Test 03 (01 Mar 2021)-img/p2_phy_s1_q01_1.png}{119.1}{0}
\qfigdef{p2_phy_s1_q02_1}{JEE Advanced (2021 Series) - Test 03 (01 Mar 2021)-img/p2_phy_s1_q02_1.png}{94.3}{0}
\qfigdef{p2_phy_s1_q03_1}{JEE Advanced (2021 Series) - Test 03 (01 Mar 2021)-img/p2_phy_s1_q03_1.png}{105.6}{0}
\qfigdef{p2_phy_s1_q04_1}{JEE Advanced (2021 Series) - Test 03 (01 Mar 2021)-img/p2_phy_s1_q04_1.png}{190.3}{0}
\qfigdef{p2_phy_s1_q08_1}{JEE Advanced (2021 Series) - Test 03 (01 Mar 2021)-img/p2_phy_s1_q08_1.png}{109.7}{0}
\qfigdef{p2_phy_s2_q09_0}{JEE Advanced (2021 Series) - Test 03 (01 Mar 2021)-img/p2_phy_s2_q09_0.png}{89.3}{0}
\qfigdef{p2_phy_s2_q09_1}{JEE Advanced (2021 Series) - Test 03 (01 Mar 2021)-img/p2_phy_s2_q09_1.png}{80.5}{0}
\qfigdef{p2_phy_s2_q10_1}{JEE Advanced (2021 Series) - Test 03 (01 Mar 2021)-img/p2_phy_s2_q10_1.png}{80.5}{0}
\qfigdef{p2_phy_s2_q11_1}{JEE Advanced (2021 Series) - Test 03 (01 Mar 2021)-img/p2_phy_s2_q11_1.png}{79.5}{0}
\qfigdef{p2_phy_s2_q12_1}{JEE Advanced (2021 Series) - Test 03 (01 Mar 2021)-img/p2_phy_s2_q12_1.png}{79.5}{0}
\qfigdef{p2_phy_s3_q02_1}{JEE Advanced (2021 Series) - Test 03 (01 Mar 2021)-img/p2_phy_s3_q02_1.png}{119.4}{0}
\qfigdef{p2_phy_s3_q03_1}{JEE Advanced (2021 Series) - Test 03 (01 Mar 2021)-img/p2_phy_s3_q03_1.png}{137.0}{0}
\qfigdef{p2_phy_s3_q04_1}{JEE Advanced (2021 Series) - Test 03 (01 Mar 2021)-img/p2_phy_s3_q04_1.png}{82.3}{0}
\qfigdef{p2_phy_s3_q06_1}{JEE Advanced (2021 Series) - Test 03 (01 Mar 2021)-img/p2_phy_s3_q06_1.png}{123.7}{0}
\qfigdef{p2_chem_s1_q04_A}{JEE Advanced (2021 Series) - Test 03 (01 Mar 2021)-img/p2_chem_s1_q04_A.png}{87.7}{0}
\qfigdef{p2_chem_s1_q04_B}{JEE Advanced (2021 Series) - Test 03 (01 Mar 2021)-img/p2_chem_s1_q04_B.png}{120.9}{0}
\qfigdef{p2_chem_s1_q04_C}{JEE Advanced (2021 Series) - Test 03 (01 Mar 2021)-img/p2_chem_s1_q04_C.png}{134.4}{0}
\qfigdef{p2_chem_s1_q04_D}{JEE Advanced (2021 Series) - Test 03 (01 Mar 2021)-img/p2_chem_s1_q04_D.png}{58.9}{0}
\qfigdef{p2_chem_s1_q06_A}{JEE Advanced (2021 Series) - Test 03 (01 Mar 2021)-img/p2_chem_s1_q06_A.png}{137.6}{0}
\qfigdef{p2_chem_s1_q06_C}{JEE Advanced (2021 Series) - Test 03 (01 Mar 2021)-img/p2_chem_s1_q06_C.png}{137.9}{0}
\qfigdef{p2_chem_s2_q09_1}{JEE Advanced (2021 Series) - Test 03 (01 Mar 2021)-img/p2_chem_s2_q09_1.png}{65.3}{0}
\qfigdef{p2_chem_s2_q09_2}{JEE Advanced (2021 Series) - Test 03 (01 Mar 2021)-img/p2_chem_s2_q09_2.png}{30.5}{0}
\qfigdef{p2_chem_s2_q09_3}{JEE Advanced (2021 Series) - Test 03 (01 Mar 2021)-img/p2_chem_s2_q09_3.png}{102.7}{0}
\qfigdef{p2_chem_s2_q09_4}{JEE Advanced (2021 Series) - Test 03 (01 Mar 2021)-img/p2_chem_s2_q09_4.png}{65.3}{0}
\qfigdef{p2_chem_s2_q09_5}{JEE Advanced (2021 Series) - Test 03 (01 Mar 2021)-img/p2_chem_s2_q09_5.png}{143.4}{0}
\qfigdef{p2_chem_s2_q09_6}{JEE Advanced (2021 Series) - Test 03 (01 Mar 2021)-img/p2_chem_s2_q09_6.png}{135.3}{0}
\qfigdef{p2_chem_s2_q10_1}{JEE Advanced (2021 Series) - Test 03 (01 Mar 2021)-img/p2_chem_s2_q10_1.png}{65.3}{0}
\qfigdef{p2_chem_s2_q10_2}{JEE Advanced (2021 Series) - Test 03 (01 Mar 2021)-img/p2_chem_s2_q10_2.png}{30.5}{0}
\qfigdef{p2_chem_s2_q10_3}{JEE Advanced (2021 Series) - Test 03 (01 Mar 2021)-img/p2_chem_s2_q10_3.png}{102.7}{0}
\qfigdef{p2_chem_s2_q10_4}{JEE Advanced (2021 Series) - Test 03 (01 Mar 2021)-img/p2_chem_s2_q10_4.png}{65.3}{0}
\qfigdef{p2_chem_s2_q10_5}{JEE Advanced (2021 Series) - Test 03 (01 Mar 2021)-img/p2_chem_s2_q10_5.png}{143.5}{0}
\qfigdef{p2_chem_s2_q10_6}{JEE Advanced (2021 Series) - Test 03 (01 Mar 2021)-img/p2_chem_s2_q10_6.png}{135.3}{0}
\qfigdef{p2_chem_s3_q01_1}{JEE Advanced (2021 Series) - Test 03 (01 Mar 2021)-img/p2_chem_s3_q01_1.png}{88.5}{0}

\begin{document}
\begin{center}{\Large\bfseries JEE Advanced (2021 Series) --- Test 03 (01 Mar 2021)}\end{center}
\begin{multicols}{2}
\papersec{1}
\partsec{1}{PHYSICS}
\sect{Section-I(i)}
\qn{1} A rod $AB$ of length $L$ slides in the $xy$-plane. If an instant angle $\theta$ with the vertical, the angular velocity of the rod will be at that instant :
\qfig{p1_phy_s1_q01_1}
\begin{qoptions}
\opt{A}{Dependent upon the length of the rod, the angle $\theta$ and also on the linear velocity of the end $A$ of the rod at that instant.}
\opt{B}{Independent of the length of the rod but will depend on the angle $\theta$.}
\opt{C}{Independent of $\theta$ but will depend on the length of the rod and linear velocity of the end $A$ at that instant.}
\opt{D}{Directly proportional to the length of the rod and the linear velocity of the end $A$ at that instant.}
\end{qoptions}
\qn{2} In two experiments with a continuous flow calorimeter to determine the specific heat capacity of a liquid, an input power of $30$ W produces a rise of $10$ K in the liquid. When the power was tripled, the same temperature rise was achieved by making the rate of flow of liquid four times faster. The power lost to the surrounding in each case is :
\begin{qoptions}
\opt{A}{$20$ W}
\opt{B}{$30$ W}
\opt{C}{$40$ W}
\opt{D}{$10$ W}
\end{qoptions}
\qn{3} $+Q$ charge is placed at midpoint of $AB$ as shown in figure. Then electric flux through cube due to $+Q$ charge is :
\qfig{p1_phy_s1_q03_1}
\begin{qoptions}
\opt{A}{$\dfrac{+Q}{\epsilon_0}$}
\opt{B}{$\dfrac{+Q}{2\epsilon_0}$}
\opt{C}{$\dfrac{+Q}{4\epsilon_0}$}
\opt{D}{$\dfrac{+Q}{8\epsilon_0}$}
\end{qoptions}
\qn{4} In a double slit experiment, instead of taking slits of equal widths, one slit is made twice as wide as the other. Then in the interference pattern
\begin{qoptions}
\opt{A}{the intensities of both the maxima and the minima increases.}
\opt{B}{the intensity of both the maxima and minima decreases.}
\opt{C}{the intensity of maxima decreases and that of the minima increases.}
\opt{D}{the intensity of maxima increases and that of the minima decreases.}
\end{qoptions}
\sect{Section-I(ii)}
\qn{5} An equiconvex lens of refractive index $n_2$ is placed such that the refractive index of the surrounding media is as shown. Then the lens :-
\qfig{p1_phy_s2_q05_1}
\begin{qoptions}
\opt{A}{must be diverging if $n_2$ is more than the arithmetic mean of $n_1$ and $n_3$.}
\opt{B}{must be converging if $n_2$ is greater than the arithmetic mean of $n_1$ and $n_3$.}
\opt{C}{must be diverging if $n_2$ is less than the arithmetic mean of $n_1$ and $n_3$.}
\opt{D}{will neither be diverging nor converging if $n_2$ is equal to arithmetic mean of $n_1$ and $n_3$.}
\end{qoptions}
\qn{6} The two balls $A$ and $B$ as shown in figure are of masses $m$ and $2m$, respectively. Ball $A$ moves with velocity $u$ toward right while $B$ is at rest. The wall at extreme right is fixed. Coefficient of restitution for collision between two balls is $\dfrac12$ and between ball and wall is $1$. Then speed of $A$ and $B$ after all possible collisions are :-
\qfig{p1_phy_s2_q06_1}
\begin{qoptions}
\opt{A}{$v_A=\dfrac u2$}
\opt{B}{$v_A=\dfrac u4$}
\opt{C}{$v_B=\dfrac u8$}
\opt{D}{$v_B=\dfrac u4$}
\end{qoptions}
\qn{7} A battery of emf $E_0=5$ V and internal resistance $5\,\Omega$ is connected across a long uniform wire $AB$ of length $1$ m and resistance per unit length $5\,\Omega\,\mathrm{m}^{-1}$. Two cells of $E_1=1$ V and $E_2=2$ V are connected as shown in figure.
\qfig{p1_phy_s2_q07_1}
\begin{qoptions}
\opt{A}{The null point is between $A$ and $B$.}
\opt{B}{If jockey is touched to point $B$, the current in the galvanometer will be going towards $B$.}
\opt{C}{When jockey is connected to point $A$, no current is flowing through $1$ V battery.}
\opt{D}{The null point is point $A$.}
\end{qoptions}
\qn{8} A car is chasing truck as shown in figure. If car sounds horn, then choose the \textbf{CORRECT} options.
\begin{center}
\qfigi{p1_phy_s2_q08_2}\hspace{1em}\qfigi{p1_phy_s2_q08_1}
\end{center}
\begin{qoptions}
\opt{A}{Apparent frequency of sound for truck driver is more than actual frequency of sound.}
\opt{B}{Apparent wavelength of sound for truck driver is more than actual wavelength of sound.}
\opt{C}{Apparent wavelength of sound for truck driver is less than actual wavelength of sound.}
\opt{D}{Speed of sound observed by truck driver is lesser than actual speed of sound in still air.}
\end{qoptions}
\qn{9} A black body of temperature $T$ is inside the chamber of temperature $T_0$ initially. Sun rays are allowed to fall from a hole from the top of chamber. If the temperature of black body ($T$) and chamber ($T_0$) remains constant, then :-
\qfig{p1_phy_s2_q09_1}
\begin{qoptions}
\opt{A}{Black body will absorb more radiation than before sun rays fall on it.}
\opt{B}{Black body will absorb less radiation than before sun rays fall on it.}
\opt{C}{Black body emits more energy than energy absorbed by it after sun rays fall on it.}
\opt{D}{Black body emits energy equal to energy absorbed by it after sun rays fall on it.}
\end{qoptions}
\qn{10} Two long thin, parallel conductors carrying equal currents in the same direction are fixed parallel to the $x$-axis, one passing through $y=a$ and the other through $y=-a$. The resultant magnetic field due to the two conductors at any point is $B$. Which of the following are \textbf{CORRECT}?
\qfig{p1_phy_s2_q10_1}
\begin{qoptions}
\opt{A}{$B=0$ for all points on the $x$-axis.}
\opt{B}{At all points on the $y$-axis, excluding the origin, $B$ has only a $z$-component.}
\opt{C}{At all points on the $z$-axis, excluding the origin, $B$ has only a $y$-component.}
\opt{D}{$B$ cannot have an $x$-component at any point in space.}
\end{qoptions}
\qn{11} When a hydrogen atom is excited from ground state to first excited state then
\begin{qoptions}
\opt{A}{Its kinetic energy increases by $10.2$ eV.}
\opt{B}{Its kinetic energy decreases by $10.2$ eV.}
\opt{C}{Its potential energy increases by $20.4$ eV.}
\opt{D}{Its angular momentum decreases by $1.05\times10^{-34}$ J-s.}
\end{qoptions}
\qn{12} An alternating voltage source is applied to series $RLC$ circuit as shown in the figure. Which of the following(s) is/are possible for the circuit at some instant of time?
\qfig{p1_phy_s2_q12_1}
\begin{qoptions}
\opt{A}{Voltage across $AD$ is zero and current through it is non-zero.}
\opt{B}{Voltage across $AD$ is zero and current through it is also zero.}
\opt{C}{Voltage across $AC$ is zero and current through it is non-zero.}
\opt{D}{Voltage across $BD$ is zero and current through it is non-zero.}
\end{qoptions}
\qfullbegin
\sect{Section-II}
\qn{1} The jaws of the Vernier callipers shown in figure are in contact with each other. Then the zero error (in mm) of this Vernier callipers is :
\qfig{p1_phy_s3_q01_1}
\qfullend
\qn{2} The minimum and maximum distances of a satellite from the centre of the earth are $2R$ and $4R$ respectively, where $R$ is the radius of earth and $M$ is the mass of the earth. The minimum speed of the satellite is $\sqrt{\dfrac{4GM}{xR}}$ then value of $x$ is :-
\qn{3} A nucleus with mass number $224$ initially at rest emits an $\alpha$ particle. If the $Q$ value of the reaction is $8$ MeV and only one photon of energy $2.00$ MeV is subsequently emitted after the emission of the $\alpha$ particle, the kinetic energy (in MeV) of the $\alpha$ particle is :
\qn{4} A positive charge particle enters and comes out from a uniform magnetic field which exists with direction inside the paper in a finite space of width $120$ cm as shown in the figure. Find the radius of the circular path of charge particle (in meter) when it is inside the magnetic field.
\qfig{p1_phy_s3_q04_1}
\qn{5} The ratio of translational and rotational kinetic energies for $1$ mole gas at $200$ K temperature is $3:2$. The internal energy of one mole gas at that temperature is $x\times10^2$ J. Then value of $x$ is $[R=8.3$ J/mol-K$]$
\qn{6} In the given figure $AB$ is a pipe fixed to the ground at an inclination of $37^\circ$. A ball is projected from point $O$ at a speed of $u=20$ m/s at angle of $53^\circ$ to the horizontal and it smoothly enter the pipe with its velocity parallel to the axis of the pipe. Then length of pipe (in metre) is :\\
$[$Take $g=10$ ms$^{-2}]$
\qfig{p1_phy_s3_q06_1}
\partsec{2}{CHEMISTRY}
\sect{Section-I(i)}
\qn{1} During the electrolysis of $0.1$ M CuSO$_4$ solution using copper electrodes, a depletion of [Cu$^{++}$] occurs near the cathode with a corresponding excess near the anode, owning to inefficient stirring of the solution. If the local concentrations of [Cu$^{++}$] near the anode and cathode are respectively $0.12$ M and $0.08$ M, calculate the back e.m.f. developed. Temperature $=298$ K\\
$[\log2=0.301,\ \log3=0.477]$
\begin{qoptions}
\opt{A}{$22$ mV}
\opt{B}{$5.2$ mV}
\opt{C}{$29$ mV}
\opt{D}{$59$ mV}
\end{qoptions}
\qn{2} Incorrect statement out of the following is :
\begin{qoptions}
\opt{A}{Beilstein Test is used as qualitative test for Organic halides}
\opt{B}{Xanthoproteic test is used to detect presence of amino acid carrying aromatic groups (Benzene ring)}
\opt{C}{Barfoed Test is used to detect the presence of Monosaccharides.}
\opt{D}{NaHSO$_3$ is used to identify C$_2$H$_5$--$\overset{\overset{\mathrm O}{\|}}{\mathrm C}$--C$_2$H$_5$.}
\end{qoptions}
\qn{3} Select the \textbf{CORRECT} statement(s) / order(s) :
\begin{qoptions}
\opt{A}{A d-orbital can accommodate $10$ electrons}
\opt{B}{\resizebox{0.96\linewidth}{!}{$\underset{[\mathrm{He}]2s^22p^63s^1}{\mathrm{Na}}\xrightarrow{\mathrm{I.E.}_1}\underset{[\mathrm{He}]2s^22p^6}{\mathrm{Na}^+}\xrightarrow{\mathrm{I.E.}_2}\underset{[\mathrm{He}]2s^22p^5}{\mathrm{Na}^{2+}}\xrightarrow{\mathrm{I.E.}_3}\underset{[\mathrm{He}]2s^22p^4}{\mathrm{Na}^{3+}}$},\\ order of successive I.E. is I.E.$_1<$ I.E.$_2>$ I.E.$_3$}
\opt{C}{Number of unpaired electrons in Co$^{2+}$ cation $>$ Number of unpaired electrons in Co$^{3+}$ cation}
\opt{D}{First ionisation energy of Zn is greater than that of Cu}
\end{qoptions}
\qn{4} The \textbf{INCORRECT} order is :
\begin{qoptions}
\opt{A}{Covalent character : MgCl$_2>$ CaCl$_2>$ SrCl$_2>$ BaCl$_2$}
\opt{B}{Thermal stability : PbF$_4>$ PbCl$_4>$ PbBr$_4>$ PbI$_4$}
\opt{C}{Melting point : KF $>$ KCl $>$ KBr $>$ KI}
\opt{D}{Boiling point : He $>$ Ne $>$ Ar $>$ Kr $>$ Xe}
\end{qoptions}
\sect{Section-I(ii)}
\qn{5} Which of the following reactions will produce benzene?
\begin{qoptions}
\opt{A}{\qfigi{p1_chem_s2_q05_1}\ $\xrightarrow{\text{Zn}}$}
\opt{B}{$3\mathrm{HC}{\equiv}\mathrm{CH}\ \xrightarrow[\Delta]{\text{Fe or Cu}}$}
\opt{C}{\qfigi{p1_chem_s2_q05_2}\ $\xrightarrow{\text{Sn + HCl}}$}
\opt{D}{\qfigi{p1_chem_s2_q05_D}\ $\xrightarrow{\text{H}_3\text{PO}_2}$}
\end{qoptions}
\qn{6} Incorrectly match options out of the following is/are.
\par\medskip\noindent\textbf{List-I}
\begin{qoptions}
\opt{A}{Maltose}
\opt{B}{Amylopectin}
\opt{C}{4 Hydroxy proline \qfigi{p1_chem_s2_q06_2}}
\opt{D}{Thymine (Nitrogenous Base)}
\end{qoptions}
\noindent\textbf{List-II}
\begin{qoptions}
\opt{P}{\qfigi{p1_chem_s2_q06_1}}
\opt{Q}{Presence of $\beta$-1,4-glycosidic and 1,6-glycosidic linkage}
\opt{R}{Will give $-$ve test with NaNO$_2$ + HCl}
\opt{S}{Presence of imidazole ring}
\end{qoptions}
\qn{7} A reaction can take place by two paths. $k_1$ and $k_2$ are rate constants for the two paths and $E_1$ and $E_2$ are their respective activation energies.\\
At temperature $T_a$ : $k_1>k_2$, $E_1<E_2$.\\
if temperature is raised to $T_b$, the rate constants change to $k_1'$ and $k_2'$. Which relation is not correct between $k_1$ and $k_2$, $k_1'$ and $k_2'$ (considering activation energy does not change with temperature).
\begin{qoptions}
\opt{A}{$\dfrac{k_1'}{k_1}>\dfrac{k_2'}{k_2}$}
\opt{B}{$\dfrac{k_1'}{k_1}=\dfrac{k_2'}{k_2}$}
\opt{C}{$\dfrac{k_1'}{k_1}<\dfrac{k_2'}{k_2}$}
\opt{D}{$\dfrac{k_1'}{k_2'}>\dfrac{k_1}{k_2}$}
\end{qoptions}
\qn{8} The variation of pH during the titration of $0.5$ N Na$_2$CO$_3$ with $0.5$ N HCl is shown in the given graph. The following table indicates the colour and pH (ranges of different indicators):
\begin{center}\scriptsize\setlength{\tabcolsep}{2pt}
\begin{tabular}{|>{\raggedright\arraybackslash}p{0.24\linewidth}|>{\raggedright\arraybackslash}p{0.27\linewidth}|>{\raggedright\arraybackslash}p{0.17\linewidth}|>{\raggedright\arraybackslash}p{0.17\linewidth}|}
\hline
\textbf{Indicator} & \textbf{Range of Colour change} & \textbf{Colour in acid} & \textbf{Colour in base}\\
\hline
Thymol blue & 1.2 to 2.8 & Red & Yellow\\
Bromocresol red & 4.2 to 6.3 & Red & Yellow\\
Bromothymol blue & 6.0 to 7.6 & Yellow & Blue\\
Cresolphthalein & 8.2 to 9.8 & Colourless & Red\\
\hline
\end{tabular}
\end{center}
\qfig{p1_chem_s2_q08_1}
Based on the graph and the table, which of the following statements are true ?
\begin{qoptions}
\opt{A}{The first equivalence point can be detected by cresolphthalein.}
\opt{B}{The complete neutralisation can be detected by bromothymol blue.}
\opt{C}{The second equivalence point an be detected by bromocresol red.}
\opt{D}{The volume of HCl required for the first equivalence point is half the volume of HCl required to reach upto second equivalence point}
\end{qoptions}
\qn{9} Xenon tetrafluoride, XeF$_4$ is/are :
\begin{qoptions}
\opt{A}{tetrahedral and act as a fluoride donor with SbF$_5$}
\opt{B}{square planar and acts as a fluoride donor with PF$_5$}
\opt{C}{square planar and acts as fluoride donor with NaF}
\opt{D}{see-saw shape and acts as a fluoride donor with AsF$_5$}
\end{qoptions}
\qn{10} Which of the following statements are \textbf{CORRECT}?
\begin{qoptions}
\opt{A}{S--S bond is present in H$_2$S$_2$O$_6$}
\opt{B}{In peroxomonosulphuric acid (H$_2$SO$_5$) sulphur is in $+6$ oxidation state}
\opt{C}{Iron powder along with Al$_2$O$_3$ and K$_2$O is used as a catalyst in the preparation of NH$_3$ by Haber's process}
\opt{D}{Changes in enthalpy is positive for the preparation of SO$_3$ by catalytic oxidation of SO$_2$}
\end{qoptions}
\qn{11} Which of the following orders are \textbf{CORRECT} ?
\begin{qoptions}
\opt{A}{TiH$_2<$ BeH$_2<$ CaH$_2$ - Electrical conductance}
\opt{B}{LiH $<$ NaH $<$ KH - Ionic character}
\opt{C}{F--F $<$ H--H $<$ D--D - Bond dissociation enthalpy}
\opt{D}{H$_2$O $<$ MgH$_2<$ NaH - Reducing character}
\end{qoptions}
\qn{12} Correctly matched with their products out of the following is/are.
\begin{qoptions}
\opt{A}{CH$_3$--CH$=$CH$_2$ $+$ Br$_2$ $\xrightarrow{\text{CCl}_4}$ CH$_3$--CHBr--CH$_2$Br\\ (Ionic addition involving cyclo bromonium ion intermediate)}
\opt{B}{CH$_3$--CH$_2$--CH$=$O $\xrightarrow{\text{H}_2\text{N--NH}_2}$\\ Nucleophilic addition to C $=$ O system}
\opt{C}{CH$_3$--CH$=$CH$_2$ $\xrightarrow[\text{Peroxide}]{\text{NBS}}$ BrCH$_2$--CH$=$CH$_2$\\ (Free Radical Allylic substitution)}
\opt{D}{\qfigi{p1_chem_s2_q12_D} $+$ Br$_2$ $\xrightarrow[\Delta]{\text{CCl}_4}$ \qfigi{p1_chem_s2_q12_3}}
\end{qoptions}
\sect{Section-II}
\qn{1} Calculate the pH at which the following conversion (reaction) will be at equilibrium in basic medium $\mathrm{I_2(s)}\rightleftharpoons\mathrm{I^-(aq)}+\mathrm{IO_3^-(aq)}$ when the equilibrium concentration at 300 K are $(\mathrm{I^-})=0.10$ M and $[\mathrm{IO_3^-}]=0.10$ M.\\
$\{$Given that $\Delta G^0_f(\mathrm{I^-,aq})=-50$ kJ/mole, $\Delta G^0_f(\mathrm{IO_3^-,aq})=-123.5$ kJ/mole, $\Delta G^0_f(\mathrm{H_2O},\ell)=-233$ kJ/mole, $\Delta G^0_f(\mathrm{OH^-,aq})=-150$ kJ/mol, ideal gas constant $=R=\dfrac{25}{3}$ J mole$^{-1}$K$^{-1}$, $\log_e10=2.3\}$
\qn{2} How many of the following statements are correct ?
\begin{qoptions}
\opt{1}{In CCP distance between two nearest octahedral voids is $2R$}
\opt{2}{In CCP distance between nearest tethrahedral and octahedral voids is $2\sqrt3R$}
\opt{3}{In CCP distance between two nearest voids is $\sqrt2R$}
\opt{4}{In CCP distance between next nearest tetrahedral void and farthest octahedral voids is $\sqrt5R$}
\opt{5}{In HCP along every edge, centres of two tetrahedral voids and one octahedral void is located}
\opt{6}{In HCP, if a plane is passed parallel to top face but at $\dfrac{3h}{4}$, from bottom it pass through centres of octahedral voids. ($h$ is the height of ionsmatic unit cell)}
\opt{7}{In HCP, top plane pass through centres of both tetrahedral and octahedral voids.}
\end{qoptions}
\qn{3} From the following reaction.
\begin{qoptions}
\opt{a}{$\mathrm{CH_3{-}\overset{\overset{O}{\|}}{C}{-}H}+x\mathrm{Cl_2}\xrightarrow{\text{NaOH}}\mathrm{Cl_3CH}+\mathrm{H{-}\overset{\overset{O}{\|}}{C}{-}ONa}$ (Haloform reaction) (`$x$' is minimum possible integer)}
\opt{b}{$y\ \mathrm{H{-}\overset{\overset{O}{\|}}{C}{-}H}+z\mathrm{NH_3}\longrightarrow$ Urotropine (Use to treat chronic urinary track infections)}
\end{qoptions}
Find out $\dfrac{x+y}{z}$.
\qn{4} How many of the following species/compounds are more acidic than carbonic acid?
\begin{center}
(a) Squaric acid \qfigi{p1_chem_s3_q04_1}\par\medskip
(b) \qfigi{p1_chem_s3_q04_2}\quad (c) \qfigi{p1_chem_s3_q04_3}\par\medskip
(d) \qfigi{p1_chem_s3_q04_4}\quad (e) \qfigi{p1_chem_s3_q04_5}\par\medskip
(f) \qfigi{p1_chem_s3_q04_6}\quad (g) HC$\equiv$C$-$CH$_3$\par\medskip
(h) MeOH\quad (i) H$_2$S
\end{center}
\qn{5} $50$ ml of water sample containing temporary hardness only, required $0.1$ ml of M/50 HCl for complete neutralisation. Calculate the temporary hardness of water in ppm.
\qn{6} If $x$ and $y$ are total number of metal electrons which are present in non-axial and axial set of d-orbitals respectively in Ni cation of $[\mathrm{Ni(DMG)_2}]$, then calculate value of $\dfrac{x^2+y^2}{x^2-y^2}$.
\partsec{3}{MATHEMATICS}
\sect{Section-I(i)}
\qn{1} A box contains $2$ fifty paise coins, $5$ twenty-five paise coins and a certain number $N\ (\ge2)$ of ten and five paise coins. Five coins are taken out of the box at random. The probability that the total value of these coins is less than one rupee and fifty paise is
\begin{qoptions}
\opt{A}{$\dfrac{10(N+2)}{{}^{N+7}C_5}$}
\opt{B}{$1-\dfrac{10(N+2)}{{}^{N+7}C_5}$}
\opt{C}{$1-\dfrac{5(N+2)}{{}^{N+7}C_5}$}
\opt{D}{None of these}
\end{qoptions}
\qn{2} If $\alpha$ and $\beta^2$ are the roots $8x^2-10x+3=0$ then the equation whose roots are $(\alpha+i\beta)^{100}$ and $(\alpha-i\beta)^{100}$ can be
\begin{qoptions}
\opt{A}{$x^2+x+1=0$}
\opt{B}{$x^2-x+1=0$}
\opt{C}{$x^3+1=0$}
\opt{D}{$2x^2+x+1=0$}
\end{qoptions}
\qn{3} If $(1+ax+bx^2)^4=a_0+a_1x+a_2x^2+\cdots+a_8x^8$ where $a,b,a_0,a_1,\dots,a_8\in R$ and $a_0+a_1+a_2\ne0$ and $\begin{vmatrix}a_0&a_1&a_2\\a_1&a_2&a_0\\a_2&a_0&a_1\end{vmatrix}=0$ the value of $5\left(\dfrac ab\right)$ is equal to
\begin{qoptions}
\opt{A}{$5$}
\opt{B}{$6$}
\opt{C}{$7$}
\opt{D}{$8$}
\end{qoptions}
\qn{4} Let $g(x)=\displaystyle\int_{1-x}^{1+x}t\,|f'(t)|\,dt$, where $f(x)$ is not a constant function in any interval $(a,b)$ and the graph of $y=f(x)$ is symmetric about the line $x=1$. Then
\begin{qoptions}
\opt{A}{$g(x)$ is increasing $x\in R$}
\opt{B}{$g(x)$ is increasing only if $x<1$}
\opt{C}{$g(x)$ is increasing if $f$ is increasing}
\opt{D}{$g(x)$ is decreasing $\forall\ x\in R$}
\end{qoptions}
\sect{Section-I(ii)}
\qn{5} The value of the definite integral $\displaystyle\int_{-\infty}^{a}\dfrac{(\sin^{-1}e^x+\sec^{-1}e^{-x})}{(\tan^{-1}e^a+\tan^{-1}e^x)(e^x+e^{-x})}\,dx$ (where $a$ is a non positive number) is
\begin{qoptions}
\opt{A}{Independent of $a$}
\opt{B}{dependent on $a$}
\opt{C}{$\dfrac\pi2\ln2$}
\opt{D}{$\dfrac\pi2\ln2\,(2\tan^{-1}e^a)$}
\end{qoptions}
\qn{6} Let $f(x)=\begin{cases}\displaystyle\int_0^x(5+|1-t|)\,dt, & \text{if } x>2\\ 5x+1, & \text{if } x\le2\end{cases}$ then the function is
\begin{qoptions}
\opt{A}{continuous at $x=2$}
\opt{B}{differentiable at $x=2$}
\opt{C}{discontinuous at $x=2$}
\opt{D}{not differentiable at $x=2$}
\end{qoptions}
\qn{7} If the solution of differential equation $(3\tan x+4\cot y-7)\sin^2y\,dx-(4\tan x+7\cot y-5)\cos^2x\,dy=0$ is $3\dfrac{\tan^2x}{2}-a\tan x+7\dfrac{\cot^2y}{2}-b\cot y+d(\cot y\tan x)=c$, ($c$ is arbitrary constant) then identify the correct options :
\begin{qoptions}
\opt{A}{$a+b+d=16$}
\opt{B}{$a+b-d=8$}
\opt{C}{$a-b+d=8$}
\opt{D}{$a\cdot d=28$}
\end{qoptions}
\qn{8} Let the unit vectors $\mathbf a$ and $\mathbf b$ be perpendicular and the unit vector $\mathbf c$ be inclined at an angle $\theta$ to both $\mathbf a$ and $\mathbf b$. If $\mathbf c=\alpha\,\mathbf a+\beta\,\mathbf b+\gamma(\mathbf a\times\mathbf b)$, ($\alpha,\beta,\gamma$ are scalars) then
\begin{qoptions}
\opt{A}{$\alpha=\beta=\cos\theta$}
\opt{B}{$\alpha=\cos\theta,\ \beta=\sin\theta$}
\opt{C}{$\gamma^2=\cos2\theta$}
\opt{D}{$\gamma^2=-\cos2\theta$}
\end{qoptions}
\qn{9} A point is moving on the circle $(x-4)^2+(y-8)^2=20$, gets broken away from it and moving along a tangent to the circle cuts the $x$-axis at point $(-2,0)$. The coordinates of the point on the circle at which the moving point broke away can be
\begin{qoptions}
\opt{A}{$\left(\dfrac{42}5,\dfrac{36}5\right)$}
\opt{B}{$\left(-\dfrac25,\dfrac{44}5\right)$}
\opt{C}{$(6,4)$}
\opt{D}{$(2,4)$}
\end{qoptions}
\qn{10} Tangent is drawn at any point $(x_1,y_1)$ other than the vertex, on parabola $y^2=4ax$. If tangents are drawn from any point on this tangent to the circle $x^2+y^2=a^2$ such that all the chords of contact pass through a fixed point $(x_2,y_2)$ then
\begin{qoptions}
\opt{A}{$x_1,\,a,\,x_2$ are in GP}
\opt{B}{$\dfrac{y_1}2,\,a,\,y_2$ are in GP}
\opt{C}{$-4,\dfrac{y_1}{y_2},\dfrac{x_1}{x_2}$ are in GP}
\opt{D}{$x_1x_2+y_1y_2=a^2$}
\end{qoptions}
\qn{11} For the figure given below if $\tan\alpha=\dfrac mn$, where $m$ and $n$ are relatively prime positive integers then
\qfig{p1_math_s2_q11_1}
\begin{qoptions}
\opt{A}{$m+n=463$}
\opt{B}{$n-m=122$}
\opt{C}{$2m-n=41$}
\opt{D}{$m+n=465$}
\end{qoptions}
\qn{12} From a point $P$ representing complex number $z_1$ on the curve $|z|=2$ two tangents are drawn from $P$ to the curve $|z|=1$ meeting it at $A(z_2)$ and $B(z_3)$ then
\begin{qoptions}
\opt{A}{centroid of $\Delta ABP$ will be on curve $|z|=1$}
\opt{B}{$\arg\left(\dfrac{z_2}{z_3}\right)=\pm\dfrac{5\pi}6$}
\opt{C}{$\left(\dfrac4{\bar z_1}+\dfrac1{\bar z_2}+\dfrac1{\bar z_3}\right)\left(\dfrac4{z_1}+\dfrac1{z_2}+\dfrac1{z_3}\right)=9$}
\opt{D}{orthocenter and circumcentre of $\Delta ABP$ will coincide.}
\end{qoptions}
\sect{Section-II}
\qn{1} A boy has $3$ library tickets and $8$ books of his interest in the library. Of these $8$, he does not want to borrow mathematics part-II, unless mathematics part-I is also borrowed. If number of ways he can choose the three books to be borrowed is $\lambda$, then number of perfect squares less than $\lambda$.
\qn{2} A slip of paper is given to a person $A$ who marks it either with a plus sign or a minus sign. The probability of his writing a plus sign is $1/3$. $A$ passes the slip to $B$, $B$ passes it to $C$, $C$ passes it to $D$ and finally $D$ passes it to a refree. Each of $B,C,D$ either leave it alone or change the sign before passing it. $B,C,D$ each change the sign with probability $2/3$. If the referee observes a plus sign on the slip then the probability that $A$ originally wrote a plus sign is
\qn{3} Two $n\times n$ square matrices $A$ and $B$ are said to be equivalent if there exists a non-singular matrix $P$ such that $P^{-1}BP=A$. If $A$ and $B$ are equivalent matrices such that $|A|=|\mathrm{adj}(\mathrm{adj}(Q))|$, where $Q=\begin{bmatrix}1&1&0\\-2&1&-1\\1&2&3\end{bmatrix}$, then sum of determinant values of $A$ and $B$ is equal to $a\times10^b$ (where `$a$' is smallest positive integer) then $\dfrac ba$ is equal to -
\qn{4} Let $f(x)=ax-bx^3$ has four extrema on $[-1,1]$ at each of which $|f(x)|=1$, where `$a$' and `$b$' are the positive integers then the value of $\dfrac ab$ is
\qn{5} Let $h(x)=(x^2+x+1)^{-1}$. The $36^\mathrm{th}$ derivative of $h(x)$ with respect to $x$ computed at $x=0$ is $k!$, then number of positive divisors of $k$ is equal to
\qn{6} Area enclosed within the curve $[y]=2^{[x]}$, $x\in[0,9)$ (where $[\cdot]$ denotes the greatest integer function) is equal to
\papersec{2}
\partsec{1}{PHYSICS}
\sect{Section-I(i)}
\qn{1} A block of mass $m$ is placed on another block of mass $2m$ as shown. A horizontal force $F$ acts on the block $m$ at height `$2b$' without any relative motion between blocks. If $F_{\min}$ is minimum force for topple the block, $a$ is acceleration of $2m$ and $f$ is friction between blocks then
\qfig{p2_phy_s1_q01_1}
\begin{qoptions}
\opt{A}{$F_{\min}=3mg$}
\opt{B}{$F_{\min}=mg/3$}
\opt{C}{$a=\dfrac g9$}
\opt{D}{$f=\dfrac{2mg}{9}$}
\end{qoptions}
\qn{2} When a capillary tube of radius $r$ is dipped in a liquid of density $\rho$ and surface tension $T$, then liquid rises to height $h$ above the liquid level. If angle of contact between liquid and tube is $60^\circ$, then
\qfig{p2_phy_s1_q02_1}
\begin{qoptions}
\opt{A}{$h=\dfrac{T}{\rho gr}$}
\opt{B}{Work done by surface tension $=\left(\dfrac{\pi T^2}{\rho g}\right)$}
\opt{C}{Increase in gravitational potential energy $=\dfrac{\pi T^2}{\rho g}$}
\opt{D}{Heat dissipated in process $=\dfrac{\pi T^2}{2\rho g}$}
\end{qoptions}
\qn{3} A disc of mass $m$ is connected with an ideal spring of stiffness $k$ is released from the unstretched position of the spring. The disc rolls down the inclined plane without slipping. Then
\qfig{p2_phy_s1_q03_1}
\begin{qoptions}
\opt{A}{Amplitude of oscillation is $\dfrac{mg\sin\theta}{k}$}
\opt{B}{Time period of oscillation is $2\pi\sqrt{\dfrac{3m}{2k}}$}
\opt{C}{Energy of oscillation is $\dfrac12\dfrac{m^2g^2\sin^2\theta}{k}$}
\opt{D}{Friction at the extreme position is $\dfrac{mg\sin\theta}{3}$}
\end{qoptions}
\qn{4} A part of circuit is shown in figure, then choose the \textbf{CORRECT} options.
\qfig{p2_phy_s1_q04_1}
\begin{qoptions}
\opt{A}{Current in $2\,\Omega$ is $2$A}
\opt{B}{$V_G-V_H=3$V}
\opt{C}{Power consumed by cell $B$ is $10$ watt}
\opt{D}{Terminal voltage across cell $A$ is $1$ volt}
\end{qoptions}
\qn{5} In a vernier callipers having $10$ vsd, the vernier constant is $0.1$ mm. When the jaws are closed, zero of vernier lies to the left of zero of main and $7^\mathrm{th}$ vsd coincides with a main scale division. When a cylinder is placed between the jaws the main scale reading was $7.7$ cm and vernier scale read $8$ division. Then
\begin{qoptions}
\opt{A}{Zero error is ($+\,0.7$) mm}
\opt{B}{Zero error is ($-0.3$)mm}
\opt{C}{Diameter of cylinder is $7.81$ cm}
\opt{D}{Diameter of cylinder is $7.71$ cm}
\end{qoptions}
\qn{6} A particle of mass $1$ kg moves along $x$-axis in a conservative force field. It's potential energy is given by $U(x)=2x^3-9x^2+12x$, where all quantities are in S.I. unit. Then select the \textbf{CORRECT} statement.
\begin{qoptions}
\opt{A}{Equilibrium positions are $x=1$ m and $x=2$m}
\opt{B}{At $x=0$ force is along $(-x)$ direction}
\opt{C}{Particle can perform oscillatory motion about $x=2$ m}
\opt{D}{For performing oscillatory motion, total mechanical energy should be less the $4$J}
\end{qoptions}
\qn{7} In a photoelectric tube, when the wavelength of incident radiation is halved, the stopping potential increases by $1$V. But when wavelength of incident radiation is increased by $100$ nm, the stopping potential is halved.
\begin{qoptions}
\opt{A}{The original stopping potential is $0.15$ V (approx.)}
\opt{B}{The original wavelength is $1240$ nm.}
\opt{C}{The work function is $0.85$ eV (approx.)}
\opt{D}{The stopping potential in $1^\mathrm{st}$ situation is $1.95$ V (approx.)}
\end{qoptions}
\qn{8} A rectangular coil in the plane of the page has dimension $a$ and $b$. A long wire that carries a current $I=2t$ is placed directly on the coil as shown. Choose the \textbf{CORRECT} options.
\qfig{p2_phy_s1_q08_1}
\begin{qoptions}
\opt{A}{Magnetic flux through coil is $\phi=\dfrac{\mu_0Ia}{2\pi}\ln\left(\dfrac{x}{b-x}\right)$ for $0<x<b$}
\opt{B}{$x\to b$, flux $\phi$ is maximum}
\opt{C}{$x=\dfrac b2$, flux $\phi$ is minimum}
\opt{D}{$x=\dfrac b4$, induced emf $=\dfrac{\mu_0a\ln3}{\pi}$}
\end{qoptions}
\sect{Section-I(ii)}
\paraq{9}{10}
When a point object, in a medium of refractive index $\mu_2$ is observed by an observer in medium of refractive index $\mu_1$, along the normal incidence, then apparent distance $=\mu_1\left[\dfrac x{\mu_1}+\dfrac y{\mu_2}\right]$. Hence, velocity of image of object as directly seen by observer is
\[v=\mu_1\left[\frac1{\mu_1}\cdot\frac{dx}{dt}+\frac1{\mu_2}\frac{dy}{dt}\right]\]
\qfig{p2_phy_s2_q09_0}
\qn{9} Answer the following by appropriately matching the lists based on the information given in the paragraph.\\
A bird in air is diving vertically over a tank with speed $6$ m/s and a fish in the tank is rising upward along the same line with speed $4$ m/sec. Base of tank is silvered. [Take $\mu=4/3$]
\qfig{p2_phy_s2_q09_1}
\par\noindent\textbf{List-I}
\begin{qoptions}
\opt{I}{Speed of the image of fish as seen by bird directly}
\opt{II}{Speed of image of fish formed after reflection from the mirror as seen by bird}
\opt{III}{Speed of image of bird relative to fish looking upwards}
\opt{IV}{Speed of image of bird relative to the fish looking downward in the mirror}
\end{qoptions}
\noindent\textbf{List-II}
\begin{qoptions}
\opt{P}{$3$ m/s}
\opt{Q}{$4$ m/s}
\opt{R}{$\dfrac92$ m/s}
\opt{S}{$9$ m/s}
\opt{T}{$12$ m/s}
\opt{U}{$\dfrac{17}2$ m/s}
\end{qoptions}
\begin{qoptions}
\opt{A}{I$\to$S ; II$\to$P ; III$\to$T, IV$\to$R}
\opt{B}{I$\to$S ; II$\to$P ; III$\to$T, IV$\to$Q}
\opt{C}{I$\to$U ; II$\to$P ; III$\to$S, IV$\to$Q}
\opt{D}{I$\to$U ; II$\to$P ; III$\to$T, IV$\to$R}
\end{qoptions}
\qn{10} Answer the following by appropriately matching the lists based on the information given in the paragraph.\\
A bird in air is diving vertically over a tank with speed $6$ m/s and a fish in the tank is also diving downwards along the same line with $4$ m/sec. Base of tank is silvered. [Take $\mu=4/3$]
\qfig{p2_phy_s2_q10_1}
\par\noindent\textbf{List-I}
\begin{qoptions}
\opt{I}{Speed of the image of fish as seen by bird directly}
\opt{II}{Speed of image of fish formed after reflection from the mirror as seen by bird}
\opt{III}{Speed of image of bird relative to fish looking upwards}
\opt{IV}{Speed of image of bird relative to the fish looking downward in the mirror}
\end{qoptions}
\noindent\textbf{List-II}
\begin{qoptions}
\opt{P}{$3$ m/s}
\opt{Q}{$4$ m/s}
\opt{R}{$\dfrac92$ m/s}
\opt{S}{$9$ m/s}
\opt{T}{$12$ m/s}
\opt{U}{$\dfrac{17}2$ m/s}
\end{qoptions}
\begin{qoptions}
\opt{A}{I$\to$S ; II$\to$P ; III$\to$T, IV$\to$R}
\opt{B}{I$\to$P ; II$\to$S ; III$\to$Q, IV$\to$U}
\opt{C}{I$\to$P ; II$\to$S ; III$\to$R, IV$\to$T}
\opt{D}{I$\to$P ; II$\to$S ; III$\to$Q, IV$\to$T}
\end{qoptions}
\paraq{11}{12}
From first law of thermodynamics, if some quantity of heat ($Q$) is supplied to a system capable of doing external work ($W$), then the quantity of heat absorbed by the system is equal to the sum of the increase in internal energy ($\Delta U$) of the system and the work done by the system
\[Q=W+\Delta U\]
Also for an ideal gas $\Delta U=\dfrac f2nR\Delta T$ where $\Delta T$ : increase in temp and $f$ : degree of freedom, Work done by force $W=\int\vec F\cdot d\vec s$, work done by gas $=\int p\cdot dv$ and gravitational potential energy $=mgh$. Take $R=\dfrac{25}3$ J / mol $-$ K, $g=10$ m/s$^2$
\qn{11} Answer the following by appropriately matching the lists based on the information given in the paragraph.\\
The piston of mass $1$kg and area of cross section $1$ cm$^2$ is attached with a cylinder as shown. The piston and cylinder are thermally insulated and contain $2$ mole of an ideal monoatomic gas. If piston is free to move and temp of gas is increased by $60$ K, choose the correct match.
\qfig{p2_phy_s2_q11_1}
\par\noindent\textbf{List-I}
\begin{qoptions}
\opt{I}{Work done by gas}
\opt{II}{Increase in internal energy of gas}
\opt{III}{Heat supplied to the gas}
\opt{IV}{Increase in gravitational potential energy of piston}
\end{qoptions}
\noindent\textbf{List-II}
\begin{qoptions}
\opt{P}{$0$}
\opt{Q}{$500$ J}
\opt{R}{$1000$ J}
\opt{S}{$1500$ J}
\opt{T}{$2000$ J}
\opt{U}{$2500$ J}
\end{qoptions}
\begin{qoptions}
\opt{A}{I$\to$R; II$\to$S; III$\to$T; IV$\to$Q}
\opt{B}{I$\to$Q ; II$\to$R; III$\to$S; IV$\to$U}
\opt{C}{I$\to$R ; II$\to$S; III$\to$U ; IV$\to$Q}
\opt{D}{I$\to$R; II$\to$S; III$\to$U; IV$\to$T}
\end{qoptions}
\qn{12} Answer the following by appropriately matching the lists based on the information given in the paragraph.\\
The piston of mass $1$kg and area of cross section $1$ cm$^2$ is attached with a cylinder as shown. The piston and cylinder are thermally insulated and contain $2$ mole of an ideal monoatomic gas. If piston is fixed and temp of gas is increased by $60$ K, choose the correct match.
\qfig{p2_phy_s2_q12_1}
\par\noindent\textbf{List-I}
\begin{qoptions}
\opt{I}{Work done by gas}
\opt{II}{Increase in internal energy of gas}
\opt{III}{Heat supplied to the gas}
\opt{IV}{Increase in gravitational potential energy of piston}
\end{qoptions}
\noindent\textbf{List-II}
\begin{qoptions}
\opt{P}{$0$}
\opt{Q}{$500$ J}
\opt{R}{$1000$ J}
\opt{S}{$1500$ J}
\opt{T}{$2000$ J}
\opt{U}{$2500$ J}
\end{qoptions}
\begin{qoptions}
\opt{A}{I$\to$P; II$\to$P; III$\to$S; IV$\to$S}
\opt{B}{I$\to$P; II$\to$Q; III$\to$S; IV$\to$P}
\opt{C}{I$\to$P; II$\to$S; III$\to$P; IV$\to$S}
\opt{D}{I$\to$P; II$\to$S; III$\to$S; IV$\to$P}
\end{qoptions}
\sect{Section-II}
\qn{1} The gap between the plates of a plane parallel plate capacitor is filled with an isotropic insulator whose dielectric varies in direction perpendicular to the plates according to law $K=\left[1+\sin\left(\dfrac{\pi x}{d}\right)\right]$ where $d$ is the separation between plates and $A$ is area of each plate. If capacitance of the capacitor is $\left(\dfrac{\lambda A\,\epsilon_0}{d}\right)$ then find the value of $\lambda$.
\qn{2} An insect of mass $2$kg jumps out of a uniform rod of same mass and length $2$m, perpendicular to the rod with a velocity $11$ m/s relative to rod. What is the angular velocity (in rad/sec) of rod just after jumping? [Assume ground is smooth]
\qfig{p2_phy_s3_q02_1}
\qn{3} A thin rod of length $2a$ is placed along $y$-axis in the $x-y$ plane. The rod carries a charge density $\lambda$ as shown in figure. If point $P_1$ is located at $(0,2a)$ and $P_2$ at $(x,0)$, then find $\left(\dfrac xa\right)$ if potential at $P_1$ and $P_2$ are equal.
\qfig{p2_phy_s3_q03_1}
\qn{4} A conical flask of height $5$m and $R=2$m is completely filled with water. Find the time (in hr.) after which the total liquid will come out of flask through a small hole of radius $1$ cm made at the bottom of flask.
\qfig{p2_phy_s3_q04_1}
\qn{5} Find the mean life (in hr) of $^{55}$Co radionuclide if its activity is known to decrease by $4\%$ per hour. The decay product is non-radioactive $\left[\ln\left(\dfrac{25}{24}\right)=0.0408\right]$,
\qn{6} In the given figure, a string of linear mass density $2$ gm/cm and length $L=2$m is stretched by a force $F=kt$ where $k$ is constant and $t$ is in sec. If at time $t=0$, a pulse is generaged at the end $A$ of the string and it will reach to end $B$ in $1$ sec. Find the value of $k$ in S.I. unit.
\qfig{p2_phy_s3_q06_1}
\partsec{2}{CHEMISTRY}
\sect{Section-I(i)}
\qn{1} In an atom, maximum no. of e$^-$ having quantum no. $n=4$, $|m_l|=2$, $m_s=\pm\dfrac12$ equals to :
\begin{qoptions}
\opt{A}{$10$}
\opt{B}{Total stereoisomers of the complex type Ma$_2$b$_2$cd}
\opt{C}{Total stereoisomers of the complex type M(A--A) b$_3$c (A--A being symmetric bidentate ligand)}
\opt{D}{Total number of 5d-subshell electrons of metal in [PtCl$_4$]$^{-2}$}
\end{qoptions}
\qn{2} Which of the following statement is \textbf{INCORRECT} ?
\begin{qoptions}
\opt{A}{O$_2$ is paramagnetic , O$_3$ is also paramagnetic}
\opt{B}{HOMO of O$_2$ is gerade}
\opt{C}{B$_2$ is paramagnetic, C$_2$ is also paramagnetic}
\opt{D}{In B$_2$ and F$_2$ s-p mixing takes place}
\end{qoptions}
\qn{3} True/correct statement for an organic compound having Molecular formula C$_5$H$_6$ is/are
\begin{qoptions}
\opt{A}{Total open chain conjugated alkadiene possible is $2$}
\opt{B}{Total possible cyclic isomer with chiral centre is $2$}
\opt{C}{Total open chain stereoisomers possible is $4$}
\opt{D}{Total open chain optically active isomers possible is $2$}
\end{qoptions}
\qn{4} Reaction which involve formation of product which has $-$CO$_2$H groups.
\begin{qoptions}
\opt{A}{\qfigi{p2_chem_s1_q04_A}}
\opt{B}{\qfigi{p2_chem_s1_q04_B}}
\opt{C}{\qfigi{p2_chem_s1_q04_C}}
\opt{D}{\qfigi{p2_chem_s1_q04_D}\ $\xrightarrow[\text{(ii) H}^+]{\substack{\text{(i) NaOH + CO}_2\\\text{(High pressure and temperature)}}}$}
\end{qoptions}
\qn{5} $\Psi=Nr\,(6-Zr)e^{-Zr/3}\cos\theta$, is a proposed hydrogenic wave function, where $Z=$ atomic number, $r=$ radial distance from the nucleus, $\theta=$ azimuthal angle, $N$ is a constant. The correct statement(s) about $\Psi$ is/are :
\begin{qoptions}
\opt{A}{$\Psi=0$ in the $xy$-plane}
\opt{B}{two radial nodes are present in $\Psi$}
\opt{C}{one angular node is present in $\Psi$}
\opt{D}{$\Psi\ne0$ in the $yz$-plane}
\end{qoptions}
\qn{6} Which reaction out the following involve formation of a hydrocarbon (cyclic/acyclic) as a major product.
\begin{qoptions}
\opt{A}{\qfigi{p2_chem_s1_q06_A}}
\opt{B}{CH$_3$CH$=$O $\xrightarrow[\text{raney Ni/H}_2]{\text{CH}_2\text{SH--CH}_2\text{SH}}$}
\opt{C}{\qfigi{p2_chem_s1_q06_C}}
\opt{D}{Ph$_3$P $\xrightarrow{\substack{\text{(i) CH}_3\text{CH}_2\text{Br}\\ \text{(ii) LDA}\\ \text{(iii) HCHO}}}$}
\end{qoptions}
\qn{7} Which of the following reduction reactions are actually employed in commercial extraction of metals ?
\begin{qoptions}
\opt{A}{Fe$_2$O$_3$ $+$ 2Al $\to$ Al$_2$O$_3$ $+$ 2Fe}
\opt{B}{Cr$_2$O$_3$ $+$ 2Al $\to$ Al$_2$O$_3$ $+$ 2Cr}
\opt{C}{2Na[Au(CN)$_2$] $+$ Zn $\to$ Na$_2$[Zn(CN)$_4$] $+$ 2 Au}
\opt{D}{Cu$_2$S $+$ Pb $\to$ Cu $+$ PbS$\downarrow$}
\end{qoptions}
\qn{8} In how many of the following cases gas is more compressible than ideal gas ? $\left(R=\dfrac{1\ \text{atm L}}{12\ \text{mol K}}\right)$
\begin{qoptions}
\opt{A}{O$_2$ gas has density $=20$ g/L at $10$ atm and $300$ K temperature.}
\opt{B}{N$_2$ gas at $2$ atm and $273$ K occupy molar volume $6$ L.}
\opt{C}{Gas $A$ at its critical temperature and pressure less than critical pressure.}
\opt{D}{CH$_4$ gas at $333$ K and in low pressure region. (Boyle's temperature $T_B$ of CH$_4=60^\circ$C)}
\end{qoptions}
\sect{Section-I(ii)}
\qn{9} Answer the following by appropriately matching the lists based on the information given in the paragraph.\\
$1^\circ$ alcohol on oxidation by mild oxidising agent produces aldehyde while with strong oxidising agent carboxylic acid is obtained. On the other hand $2^\circ$ alcohol with mild oxidising agent gives ketone. $3^\circ$ alcohol is reluctant to oxidation by mild oxidising agents.
\par\noindent\textbf{List-I (Reaction)}
\begin{qoptions}
\opt{I}{\qfigi{p2_chem_s2_q09_1}\qfigi{p2_chem_s2_q09_2}}
\opt{II}{\qfigi{p2_chem_s2_q09_3}\qfigi{p2_chem_s2_q09_4}}
\opt{III}{\qfigi{p2_chem_s2_q09_5}}
\opt{IV}{\qfigi{p2_chem_s2_q09_6}}
\end{qoptions}
\noindent\textbf{List-II (Reagents)}
\begin{qoptions}
\opt{P}{PCC, CH$_2$Cl$_2$}
\opt{Q}{KMnO$_4$ / H$^+$}
\opt{R}{Collin's Reagent}
\opt{S}{Cu/300$^\circ$C}
\opt{T}{TsCl $+$ DMSO $+$ NaHCO$_3$}
\opt{U}{MnO$_2$, $\Delta$}
\end{qoptions}
Incorrectly matched among the following is :
\begin{qoptions}
\opt{A}{I $-$ S}
\opt{B}{II $-$ P}
\opt{C}{IV $-$ S}
\opt{D}{III $-$ Q}
\end{qoptions}
\qn{10} Answer the following by appropriately matching the lists based on the information given in the paragraph.\\
$1^\circ$ alcohol on oxidation by mild oxidising agent produces aldehyde while with strong oxidising agent carboxylic acid is obtained. On the other hand $2^\circ$ alcohol with mild oxidising agent gives ketone. $3^\circ$ alcohol is reluctant to oxidation by mild oxidising agents.
\par\noindent\textbf{List-I (Reaction)}
\begin{qoptions}
\opt{I}{\qfigi{p2_chem_s2_q10_1}\qfigi{p2_chem_s2_q10_2}}
\opt{II}{\qfigi{p2_chem_s2_q10_3}\qfigi{p2_chem_s2_q10_4}}
\opt{III}{\qfigi{p2_chem_s2_q10_5}}
\opt{IV}{\qfigi{p2_chem_s2_q10_6}}
\end{qoptions}
\noindent\textbf{List-II (Reagents)}
\begin{qoptions}
\opt{P}{PCC, CH$_2$Cl$_2$}
\opt{Q}{KMnO$_4$ / H$^+$}
\opt{R}{Collin's Reagent}
\opt{S}{Cu/300$^\circ$C}
\opt{T}{TsCl $+$ DMSO $+$ NaHCO$_3$}
\opt{U}{MnO$_2$}
\end{qoptions}
Correctly matched away the following is :
\begin{qoptions}
\opt{A}{II $-$ P}
\opt{B}{II $-$ Q}
\opt{C}{II $-$ S}
\opt{D}{II $-$ U}
\end{qoptions}
\qn{11} Matching the following :
\par\noindent\textbf{List-I (complex ion)}
\begin{qoptions}
\opt{I}{[Cr(NH$_3$)$_6$]$^{3+}$}
\opt{II}{[Cu(NH$_3$)$_6$]$^{2+}$}
\opt{III}{[Fe(H$_2$O)$_6$]$^{3+}$}
\opt{IV}{[IrF$_6$]$^{3-}$}
\end{qoptions}
\noindent\textbf{List-II (CFSE, hybridisation and magnetic nature)}
\begin{qoptions}
\opt{P}{$0.0\ \Delta_0$, sp$^3$d$^2$-hybridisation}
\opt{Q}{$-0.6\ \Delta_0$, sp$^3$d$^2$-hybridisation}
\opt{R}{$-1.2\ \Delta_0$, d$^2$sp$^3$-hybridisation}
\opt{S}{Diamagnetic}
\opt{T}{Paramagnetic}
\opt{U}{$0.0\ \Delta_0$, d$^2$sp$^3$ hybridisation}
\end{qoptions}
\textbf{CORRECTLY} matched among the following is :
\begin{qoptions}
\opt{A}{II $-$ P}
\opt{B}{IV $-$ S}
\opt{C}{I $-$ S}
\opt{D}{III $-$ R}
\end{qoptions}
\qn{12} Matching the following :
\par\noindent\textbf{List-I (complex ion)}
\begin{qoptions}
\opt{I}{[Cr(NH$_3$)$_6$]$^{3+}$}
\opt{II}{[Cu(NH$_3$)$_6$]$^{2+}$}
\opt{III}{[Fe(H$_2$O)$_6$]$^{3+}$}
\opt{IV}{[IrF$_6$]$^{3-}$}
\end{qoptions}
\noindent\textbf{List-II (CFSE, hybridisation and magnetic nature)}
\begin{qoptions}
\opt{P}{$0.0\ \Delta_0$, sp$^3$d$^2$-hybridisation}
\opt{Q}{$-0.6\ \Delta_0$, sp$^3$d$^2$-hybridisation}
\opt{R}{$-1.2\ \Delta_0$, d$^2$sp$^3$-hybridisation}
\opt{S}{Diamagnetic}
\opt{T}{Paramagnetic}
\opt{U}{$0.0\ \Delta_0$, d$^2$sp$^3$ hybridisation}
\end{qoptions}
\textbf{CORRECTLY} match among the following is :
\begin{qoptions}
\opt{A}{IV $-$ T}
\opt{B}{I $-$ Q}
\opt{C}{II $-$ Q}
\opt{D}{III $-$ S}
\end{qoptions}
\sect{Section-II}
\qn{1} $1$ mol of an ideal gas [$\gamma=1.5$], is taken through a series of processes as shown in following in P-V curve.
\qfig{p2_chem_s3_q01_1}
Information-I : The temperature at state `$A$' is $300$ K.\\
Information-II : The reversible isothermal expansion from $B$ to $C$ doubles the volume.\\
Information-III : The entropy change of the system from $C$ to $D$ is $4\,\ell n\,16$ cal/K.\\
How many of the following statements are correct for above information :
\begin{qoptions}
\opt{A}{The entropy change of the system from $D$ to $A$ is $-3R\ \ell n8$}
\opt{B}{the temperature at point $D$ is $2400$ K}
\opt{C}{The work done from $A$ to $B$ is $-600$ cal}
\opt{D}{The work done from $B$ to $C$ is $-300\ \ell n2$ cal}
\opt{E}{The pressure at point $A$ is $8$ times of $B$}
\opt{F}{The temperature at point $D$ is $600$ K}
\opt{G}{The pressure at point $C$ is half of pressure at point $B$.}
\end{qoptions}
\qn{2} Small spherical ball of silver metal used in jewellery having diameter $0.1$ cm, which is obtained by the electrolytic deposition. It total number of balls in jewellery is $10{,}000$, then calculate the applied amount of electricity in coulombs, which is used to make required number of balls. [Density of Ag $=10.5$, gm/cm$^3$]\\
It is assumed that $3.5\%$ electricity consumed as wastage during electrolysis [Give your answer in multiple of $10^4$] [Ag $=108$]
\qn{3} $\underset{\text{(Alkyl Halide)}}{A}+\underset{\text{(Organometallic Compound)}}{B}\xrightarrow{\text{dry ether}}\underset{(\mathrm{C_xH_y})}{C}+\mathrm{Mg}\begin{matrix}\text{Br}\\\text{Cl}\end{matrix}$
\begin{qoptions}
\opt{a}{Branched structural isomer of $A$ i.e. $A^*$ with Mol. mass $78.5$g is optically inactive.}
\opt{b}{$133$ g of $B$ when treated completely with $40$g of $D$ (acyclic Hydrocarbon) produces $1$ mole of $E$ (a Hydrocarbon)}
\opt{c}{Product $C$ give $4$ monochlorinated product on treatment with Cl$_2$/h$\nu$ (including stereo)}
\end{qoptions}
From the above information find out the Mol. mass of next successive Homologue of $C$.
\qn{4} A sample of $_{53}\mathrm{I}^{131}$, as I$^-$ ion, was administered to a patient in a carrier consisting $1.0$ mg of stable I$^-$ ion. After $4.0$ days, $60\%$ of the initial radioactivity was detected in the thyroid gland of the patient. What mass (in mg) of the stable I$^-$ ion had migrated to the thyroid gland ? Given $t_{1/2}$ to I$^{131}=8$ days
\qn{5} How many of the following statements is/are correct ?
\begin{qoptions}
\opt{1.}{Adsorbate is that substance where adsorption takes place.}
\opt{2.}{Lower the critical temperature of the gas, the more readily it will be adsorbed.}
\opt{3.}{Physical adsorption is an irreversible process.}
\opt{4.}{The amount of the gas adsorbed is directly proportional to the pressure, in all pressure region.}
\opt{5.}{By the gas masks, the poisonous gases are removed from air by charcoal, using process of adsorption.}
\opt{6.}{The water freed from all ions (cations and anions) is referred as mineral water.}
\opt{7.}{Heat of adsorption is defined as the energy liberated when $1$ gram of a gas is adsorbed on the solid surface.}
\opt{8.}{The adsorption of hydrogen on charcoal is chemical adsorption.}
\opt{9.}{Froth floatation process is an application of adsorption.}
\end{qoptions}
\qn{6} $\underset{\text{Black}}{\mathrm{X(s)}}\xrightarrow{\text{dil. HCl}}\mathrm{Y}\!\uparrow\xrightarrow[\text{NaOH}]{\mathrm{Na_2[Fe(CN)_5(NO)]}}$ Purple solution\\
Gas $Y$ has been allowed to react with following species in neutral/acidic medium :\\
(a) FeCl$_3$\quad (b) CuSO$_4$\quad (c) BaCl$_2$\quad (d) SO$_2$\quad (e) Cr$_2$O$_7^{2-}$\quad (f) CH$_3$COONa\quad (g) Hg$^{2+}$\\
Then calculate value of $(P-Q)/R$ :\\
$P$ : Number of species which undergoes redox reaction with gas $Y$.\\
$Q$ : Number of species with which gas $Y$ undergoes precipitation without going into redox reaction.\\
$R$ : Number of species with gas $Y$ produce no observable change.
\partsec{3}{MATHEMATICS}
\sect{Section-I(i)}
\qn{1} Vertices A, B and C of a tetrahedron ABCD are $(1,1,1)$, $(1,0,0)$, $(3,0,0)$ respectively. The altitude from vertex D to the opposite face ABC meets the median line through A of the $\Delta$ABC at a point E (where A \& E lie on opposite sides of BC). If the length of side AD is $4$ and volume of the tetrahedron is $\dfrac{2\sqrt2}{3}$, then the correct statement(s) is/are
\begin{qoptions}
\opt{A}{the altitude from vertex D is $2$}
\opt{B}{EB $=\sqrt6$}
\opt{C}{area$(\Delta\text{EBC})=\dfrac{2}{\sqrt3}$}
\opt{D}{vector $\hat j-\hat k$ is normal to the plane ABC}
\end{qoptions}
\qn{2} Consider $f(x)=\tan^{-1}\left(\dfrac{|x|}{\sqrt{1-x^2}}\right)+\cos^{-1}\left(\dfrac{1}{\sqrt{1+x^2}}\right)$ then
\begin{qoptions}
\opt{A}{maximum value of $f(x)$ is $\dfrac{3\pi}{4}$}
\opt{B}{number of integers in the range of $f(x)$ is $3$}
\opt{C}{number of solutions of $f(x)=2$ is $2$}
\opt{D}{if $\theta\in\left(\dfrac\pi4,\dfrac\pi2\right)$ then $f(\sin\theta)>f(\cos\theta)$}
\end{qoptions}
\qn{3} Consider a hyperbola $x^2-y^2-2x+4y-7=0$. If its pair of asymptotes meet two tangents each of slope $2$ at P, Q and R, S (in order) then
\begin{qoptions}
\opt{A}{PQ $=\dfrac{4\sqrt5}{\sqrt3}$}
\opt{B}{Latus rectum of hyperbola is $4$}
\opt{C}{area of quadrilateral PQRS is $12$ sq.units}
\opt{D}{area of quadrilateral PQRS is $16$ sq.units.}
\end{qoptions}
\qn{4} If $x,y,z$ are three distinct real numbers such that $x,12,y$ are in H.P. and $x,12,z,y$ are in increasing A.P. then
\begin{qoptions}
\opt{A}{$x+z=18$}
\opt{B}{$x+z=24$}
\opt{C}{maximum value of\\ $\sqrt{(x-3)\sin\alpha-(y-10)\cos\beta+2}$ is $4$, where $\alpha,\beta\in\mathbb R$}
\opt{D}{maximum value of $\sqrt{(x-3)\sin\alpha-(y-10)\sin\alpha}$ is $10$, where $\alpha\in\mathbb R$}
\end{qoptions}
\qn{5} Let A$(z_1)$, B$(z_2)$, C$(z_3)$ and D$(z_4)$ be vertices (in anti-clockwise order) of a trapezium (AB $\parallel$ CD) in an argand plane. Let $|z_1-z_2|=4$, $|z_3-z_4|=10$ and the diagonals AC and BD intersect at P. It is given that $\arg\left(\dfrac{z_4-z_2}{z_3-z_1}\right)=\dfrac\pi2$ and $\arg\left(\dfrac{z_3-z_2}{z_4-z_1}\right)=-\dfrac\pi4$ then
\begin{qoptions}
\opt{A}{Area of Trapezium is $\dfrac{130}{3}$ sq. unit}
\opt{B}{$|\text{CP}-\text{DP}|=\dfrac{10}{\sqrt{21}}$ unit}
\opt{C}{area of $\Delta$PCB is $\dfrac{200}{21}$ sq. unit}
\opt{D}{area of $\Delta$PCB is $\dfrac{100}{21}$ sq. unit}
\end{qoptions}
\qn{6} If in a triangle ABC, $\sum\sin\dfrac A2=\dfrac65$ and $\sum\text{I I}_1=9$ (where $\text{I}_1,\text{I}_2,\text{I}_3$ are ex-centres and I is incentre) then
\begin{qoptions}
\opt{A}{the circumradius R is $15/4$}
\opt{B}{the circumradius R is $15/8$}
\opt{C}{the value of $\sin\left(\dfrac{\pi-A}{4}\right)\cdot\sin\left(\dfrac{\pi-B}{4}\right)\cdot\sin\left(\dfrac{\pi-C}{4}\right)$ is $\dfrac1{16}$}
\opt{D}{the value of $\sin\left(\dfrac{\pi-A}{4}\right)\cdot\sin\left(\dfrac{\pi-B}{4}\right)\cdot\sin\left(\dfrac{\pi-C}{4}\right)$ is $\dfrac1{20}$}
\end{qoptions}
\qn{7} Two sides AB and BC of a $\Delta$ABC be represented by $x-7y+3=0$ and $x+y+3=0$ respectively. Incentre I $(-3-3\alpha,\alpha),\alpha\in\mathbb R$ of the triangle lies on the altitude $x-y-1=0$ through A then
\begin{qoptions}
\opt{A}{The equation of a line representing side AC is $7x-y=11$}
\opt{B}{The equation of a line representing side AC is $7x+y=1$}
\opt{C}{The area of $\Delta$ABC is $\dfrac{32}{3}$ sq. unit}
\opt{D}{The area of $\Delta$ABC is $\dfrac{32}{5}$ sq.unit}
\end{qoptions}
\qn{8} If $\vec x\times\vec y=\vec a$, $\vec y\times\vec z=\vec b$, $\vec x\cdot\vec b=\gamma$, $\vec x\cdot\vec y=1$ \& $\vec y\cdot\vec z=1$ then $\left(|\vec y|\gamma\neq0\right)$
\begin{qoptions}
\opt{A}{$\vec y=\dfrac{\vec a\times\vec b}{\gamma}$}
\opt{B}{$\vec x=\dfrac1{y^2}\left[\vec y-\vec a\times\vec y\right]$}
\opt{C}{$\vec z=\dfrac1{y^2}\left[\vec y+\vec b\times\vec y\right]$}
\opt{D}{$\vec x=\dfrac1{y^2}\left[\vec y+\vec a\times\vec y\right]$}
\end{qoptions}
\sect{Section-I(ii)}
\paraq{9}{10}
Let $A=\{1,3,5,7\}$ and $B=\{2,4,6,8\}$ and $f:A\to B$ is a function defined from set A to set B.
\par\noindent\textbf{List-I}
\begin{qoptions}
\opt{I}{Number of functions, that satisfy condition $i+f(i)<10$; where $i=1,3,5,7$; is}
\opt{II}{Number of functions, that satisfy $f(i)\neq1+i$, where $i=1,3,5,7$; is}
\opt{III}{Number of one-one and onto function, that satisfy condition $f(i)\neq1+i$, where $i=1,3,5,7$; is}
\opt{IV}{If number of many one functions from $f:A\to B$ is $k$, then $k$ is divisible by}
\end{qoptions}
\noindent\textbf{List-II}
\begin{qoptions}
\opt{P}{$24$}
\opt{Q}{$17$}
\opt{R}{$9$}
\opt{S}{$29$}
\opt{T}{$81$}
\opt{U}{$14$}
\end{qoptions}
\qn{9} Answer the following by appropriately matching the lists based on the information given in the paragraph.\\
Which of the following is correct combination ?
\begin{qoptions}
\opt{A}{(I), (S)}
\opt{B}{(II), (P)}
\opt{C}{(III) (R)}
\opt{D}{(IV) (Q)}
\end{qoptions}
\qn{10} Answer the following by appropriately matching the lists based on the information given in the paragraph.\\
Which of the following is correct combination ?
\begin{qoptions}
\opt{A}{(I) (T)}
\opt{B}{(II) (U)}
\opt{C}{(III) (Q)}
\opt{D}{(IV) (S)}
\end{qoptions}
\paraq{11}{12}
Consider, $f(x)=\dfrac{x^2}{1+x^2}$ and $g(x)=px$, $p\in\mathbb R$.
\par\noindent\textbf{List-I}
\begin{qoptions}
\opt{I}{If $f(x)=g(|x|)$ has exactly three distinct solutions then value of $2p$ is}
\opt{II}{If $f(x)=g(|x|)$ has exactly five distinct solutions then no. of integers in range of $p$ is}
\opt{III}{Number of distinct tangents that can be drawn to the curve $y=f(x)$ from the origin is}
\opt{IV}{Number of points on curve $\dfrac{g(x)}{f(x)}$, $(p\neq0)$ where tangent is parallel to x-axis, is}
\end{qoptions}
\noindent\textbf{List-II}
\begin{qoptions}
\opt{P}{$0$}
\opt{Q}{$5$}
\opt{R}{$1$}
\opt{S}{$2$}
\opt{T}{$4$}
\opt{U}{$3$}
\end{qoptions}
\qn{11} Answer the following by appropriately matching the lists based on the information given in the paragraph.\\
Which of the following is correct combination ?
\begin{qoptions}
\opt{A}{(I)(Q)}
\opt{B}{(II)(T)}
\opt{C}{(III)(R)}
\opt{D}{(IV)(S)}
\end{qoptions}
\qn{12} Answer the following by appropriately matching the lists based on the information given in the paragraph.\\
Which of the following is correct combination ?
\begin{qoptions}
\opt{A}{(I)(U)}
\opt{B}{(II)(P)}
\opt{C}{(III)(S)}
\opt{D}{(IV)(T)}
\end{qoptions}
\sect{Section-II}
\qn{1} If the value of $\displaystyle\sum_{k=0}^{7}\left(\dfrac{{}^7C_k}{{}^{14}C_k}\sum_{r=k}^{14}\left({}^rC_k\right)\left({}^{14}C_r\right)\right)$ is $a^b$, where $a$ and $b$ are relative prime, then $(a+b)$ is
\qn{2} If $\alpha,\beta,\gamma,\delta$ are the roots of the equation $x^4-x^3-x^2-1=0$ and consider $P(x)=x^6-x^5-x^3-x^2-x$. Then the value of $P(\alpha)+P(\beta)+P(\gamma)+P(\delta)$ is
\qn{3} If $\displaystyle\int\dfrac{\left(\dfrac23+x^2\right)dx}{3-4\sqrt3x+9x^4}=\dfrac1{6\sqrt6}\tan^{-1}\left(\dfrac{\sqrt3x+1}{\sqrt2}\right)-\dfrac1{6\sqrt3}\cdot\dfrac1{f(x)}+C$, where C is constant of integration and $x\neq\dfrac1{\sqrt3}$, $f(0)=-1$, then value of $f\left(\sqrt3\right)$ is
\qn{4} Let $f:(-1,\infty)\to(0,\infty)$ be a differentiable function satisfying $f(x)\cdot f(yf(x))=f(x+y)\ \forall\ x,y$ (for which $f$ is well defined in given domain) and $f'(0)=-1$ and $f(0)=1$ then area of the region bounded by the curve $y=\dfrac{f(x)}{x+1}$ and the coordinate axes is
\qn{5} The value of $\displaystyle\lim_{x\to\infty}\dfrac{\left(xe^{-x}+x^2e^{-2x}+x^3e^{-3x}+\dots\infty\right)}{\left(\ln(1+x)\right)^{-1}}$ is equal to
\qn{6} Suppose a function $f(x)$ has an inverse function $g(x)$. If $f(3)=4$, $f'(3)=\left(\dfrac25\right)^{1/3}$ and $f''(3)=-\dfrac47$, then $g''(4)$ is equal to
\end{multicols}
\end{document}
